\section{Homology and the finite quotients}
\label{sec:homology}

We now prove that the finite complexes $T_m$ have no cohomology in
degrees $1,\ldots,d-1$. We write $H_j(G;A)$ for group homology with
trivial coefficients in $A$. The acyclic complex $L$ first compares
the homology of $U$ and its congruence subgroups with simplicial homology.
We then use stabilization to show that the congruence subgroups inherit
the vanishing of $H_j(U;\F)$.

\subsection{An acyclic complex as a partial resolution}
\label{subsec:projective-resolution}

Finite odd-order stabilizers make the simplicial chains projective over
the group algebra. This allows the acyclicity already proved for $L$
to replace a full contractible model in the degrees we need.

\begin{lemma}\label{lem:projective-chains}
Let $h\ge0$ be an integer, and let a group $K$ act on a nonempty
simplicial complex $Z$ without inversions. Suppose every nonempty
simplex stabilizer is finite of odd order.
If $\wtH_j(Z;\F)=0$ for $0\le j\le h$, then
\[
 H_j(K;\F)\cong H_j(C_*(Z;\F)_K)\qquad(0\le j\le h).
\]
Here the subscript denotes coinvariants; this is the chain complex of
simplex orbits, with boundary induced by deleting vertices.
\end{lemma}

\begin{proof}
For every $r\ge0$, the permutation module $C_r(Z;\F)$ is a direct sum
of modules $\F[K/H]$, where $H$ runs over simplex stabilizers. Because
$|H|$ is odd, averaging over $H$ splits the augmentation
$\F[H]\to\F$. Thus $\F$ is a projective $\F[H]$-module, and induction
makes $\F[K/H]$ a projective $\F[K]$-module.
The augmented chains of $Z$ are exact through degree $h$. Retain the
chain modules through degree $h+1$, choose a projective module mapping
onto the kernel of the boundary out of $C_{h+1}(Z;\F)$, and continue
in higher degrees. This gives a
projective resolution of the trivial module $\F$ with the same terms
and differentials in all degrees needed to compute $H_j$, $j\le h$.
Taking coinvariants proves the assertion.
\end{proof}

\begin{proposition}\label{prop:homology-comparisons}
For $0<j<n-1$ and $m\ge2$,
\begin{equation}\label{eq:homology-comparisons}
 H_j(U;\F)=0,\qquad
 H_j(T_m;\F)\cong H_j(\Gamma_N(m);\F).
\end{equation}
Moreover, $T_m$ is connected.
\end{proposition}

\begin{proof}
For a nonempty type set $A\subseteq S$, the simplices of $L$ of type $A$
are the cosets of $\bigcap_{i\in A}H_i=U_{S\setminus A}$. Thus $U$ is
transitive on the faces of each type, and the orbit complex $U\backslash L$
is a single $n$-simplex. Its stabilizers are conjugates of the finite
unitriangular groups in \cref{lem:patterns}, so they have odd order.
Apply \cref{lem:projective-chains} and \cref{prop:L-acyclic} to obtain
the first assertion of \eqref{eq:homology-comparisons}. In the coinvariant
boundary each face is obtained by deleting one vertex, with coefficient
$1$; there are no stabilizer-index factors.

Reduction modulo $t^m$ preserves all the stabilizer intersections by
\cref{lem:patterns}. To check the resulting quotient, let two type-$A$
simplices have representatives $u,v\in U$ and the same reduced vertices.
Then the reduction of $u^{-1}v$ belongs to
$\bigcap_{i\in A}\overline H_i=\overline{U_{S\setminus A}}$.
Choose $h\in U_{S\setminus A}$ with this reduction. The element
$vh^{-1}u^{-1}\in\Gamma_N(m)$ carries the first simplex to the second.
Thus reduction identifies precisely the simplex orbits and gives
\begin{equation}\label{eq:quotient}
 \Gamma_N(m)\backslash L\cong T_m.
\end{equation}
The $\Gamma_N(m)$-stabilizers of nonempty simplices are trivial, since
the reduction is injective on every local group and $\Gamma_N(m)$ is
normal in $U$. The action preserves types. A second application of
\cref{lem:projective-chains} gives the remaining homology comparison.
Connectedness follows from \cref{prop:L-acyclic} and \eqref{eq:quotient}.
\end{proof}

\subsection{From pro-stability to ordinary surjectivity}
\label{subsec:pro-stability}

Our next goal is a surjection on homology at each fixed congruence level.
The available stability theorem concerns inverse systems; odd-index
transfer will supply the additional step from that theorem to an
ordinary surjection.

We first specify its ring conventions. For a commutative nonunital ring
$A$, its \emph{integer unitization} is $\Z\ltimes A=\Z\oplus A$, with
\[
 (u,a)(v,b)=(uv,ub+va+ab),\qquad 1=(1,0).
\]
The projection to $\Z$ defines
\[
 \GL_a(A)=\ker\bigl(\GL_a(\Z\ltimes A)\longrightarrow\GL_a(\Z)\bigr).
\]
A row $b\in(\Z\ltimes A)^\ell$ is \emph{$A$-unimodular} if
$b\equiv e_1\pmod A$ and there is a row $c\equiv e_1\pmod A$ with
$b\cdot c=1$. Here $e_1=(1,0,\ldots,0)$. We use the relative stable
range convention of Vaserstein \cite[Section~2, pp.~384--385]{Vas1969}:
$\sr(A)\le s$ means that, for every such row of length $\ell\ge s+1$,
there are $v_1,\ldots,v_{\ell-1}\in A$ for which
\[
 (b_1+v_1b_\ell,\ldots,b_{\ell-1}+v_{\ell-1}b_\ell)
\]
is $A$-unimodular.

An inverse system $(A_m)$ is \emph{pro-zero} if, for every index, some later
transition to that index is zero. A level map of inverse systems of
abelian groups is a \emph{pro-isomorphism} precisely when its kernel
and cokernel systems are pro-zero. The system $(A_m)$ is \emph{Tor-unital}
if $(\Tor_i^{\Z\ltimes A_m}(\Z,\Z))_m$ is pro-zero for every $i>0$.
Morrow's theorem says that the powers of any ideal in a commutative
Noetherian ring form a Tor-unital system
\cite[Theorem~0.3, arXiv~v2]{Morrow}.
For a Tor-unital inverse system of commutative nonunital rings,
Iwasa's integral homology pro-stability theorem says that, if
$r=\max\{2,\sup_m\sr(A_m)\}<\infty$, then
\begin{equation}\label{eq:iwasa}
 (H_j(\GL_a(A_m);\Z))_m\longrightarrow
 (H_j(\GL_{a+1}(A_m);\Z))_m
 \quad\text{is a pro-isomorphism for }a\ge2j+r-1.
\end{equation}
We cite the published version \cite[Theorem~5.13]{Iwasa}.

\begin{lemma}\label{lem:stabilization}
Let $h\ge0$ be an integer. If $a\ge2h+1$, the inclusion on the first
$a$ coordinates induces
surjections
\[
 H_j(\Gamma_a(m);\F)\longrightarrow H_j(\Gamma_{a+1}(m);\F)
 \qquad(0\le j\le h,\ m\ge1).
\]
\end{lemma}

\begin{proof}
Apply \eqref{eq:iwasa} to $A_m=t^m k[t]$. In this case
$\GL_a(A_m)$ is exactly the group of matrices in $\GL_a(R)$ congruent
to $I$ modulo $t^m$: the entrywise map from the unitization is injective
on $A_m$, and an inverse congruent to $I$ lifts back to the unitization.
Such a matrix has determinant in $R^*=k^*$, congruent to $1$ modulo $t^m$;
hence its determinant is $1$. Thus $\GL_a(A_m)=\Gamma_a(m)$.

The nonunital stable range of every $A_m$ is at most $2$. Here is a direct
verification of the relative row criterion. For a unimodular row
$(b_1,\ldots,b_r)$ with $r\ge3$ and $(b_1,\ldots,b_r)\equiv(1,0,\ldots,0)
\pmod{t^m}$, the polynomial $b_1$ is coprime to $t$. At each irreducible
divisor of $b_1$ at which $b_r$ is nonzero, choose a residue for $v$ avoiding
$-b_2/b_r$. The Chinese remainder theorem allows all these choices with
$v\in t^mR$. Then
$(b_1,b_2+vb_r,b_3,\ldots,b_{r-1})$ is unimodular and still has the same
residue modulo $t^m$: at a divisor where $b_r$ vanishes, original
unimodularity already makes one of $b_2,\ldots,b_{r-1}$ nonzero.
To check the unitization condition, if a row $b$
has $b\cdot c=1$ over $R$ and $b\equiv e_1\pmod{t^m}$, then
\[
 c'=e_1+(1-b_1)c
 \quad\text{satisfies}\quad b\cdot c'=1,
 \qquad c'\equiv e_1\pmod{t^m}.
\]
It therefore lifts to a unimodularity relation over $\Z\ltimes A_m$:
the difference between the lifted dot product and $1$ lies in $A_m$,
where the map to $R$ is injective.
This proves the claimed stable-range bound.

Morrow's theorem and \eqref{eq:iwasa} now give pro-isomorphisms in integral
homology for $0\le j\le h$. The natural universal coefficient sequences
\[
 0\longrightarrow H_j(G;\Z)\otimes\F
 \longrightarrow H_j(G;\F)
 \longrightarrow\Tor_1^{\Z}(H_{j-1}(G;\Z),\F)
 \longrightarrow0
\]
give pro-isomorphisms with $\F$ coefficients in the same range. For
$j\ge1$, both integral degrees $j$ and $j-1$ are covered; degree zero
is immediate. Indeed,
the tensor and Tor functors extend to inverse systems and preserve
isomorphisms in the pro-category, and these levelwise exact sequences
are exact there as well. In particular, the systems of cokernels of the
displayed stabilization maps over $\F$ are pro-zero.

For $m'\ge m$, the index $[\Gamma_a(m):\Gamma_a(m')]$ is finite and odd.
To see this, reduce modulo $t^{m'}$ and filter matrices congruent to $I$
modulo $t^m$ by successive powers of $t$; each successive quotient embeds
into an additive matrix group over $k$. Transfer therefore gives
\[
 H_j(\Gamma_a(m);\F)\xrightarrow{\operatorname{tr}}
 H_j(\Gamma_a(m');\F)\xrightarrow{\operatorname{inc}_*}
 H_j(\Gamma_a(m);\F),
 \qquad \operatorname{inc}_*\operatorname{tr}=1.
\]
Thus all homology transition maps, and consequently all cokernel
transition maps, are surjective. A pro-zero system with surjective
transition maps is zero at every level: a sufficiently late transition
to any chosen term is both zero and surjective. Applying this observation
to the cokernels gives the required ordinary surjections.
\end{proof}

\subsection{Trivial conjugation and vanishing cohomology}
\label{subsec:conjugation}

Stabilization lets us represent every homology class on any chosen set
of $N-2$ coordinates. The two remaining coordinates will make each
elementary conjugation act trivially. This is the finite-level argument
corresponding to the stability-to-conjugation principle in
\cite[Corollary~5.14]{Iwasa}.

\begin{lemma}\label{lem:trivial-action}
For $1\le j\le d-1$, conjugation by $\SL_N(R)$ acts trivially on
$H_j(\Gamma_N(m);\F)$, for every $m\ge1$.
\end{lemma}

\begin{proof}
Our choice of $N$ gives $N-2=2(d-1)+1$. By two applications of
\cref{lem:stabilization}, the subgroup $\Gamma_{N-2}(m)$ on the first
$N-2$ coordinates surjects onto $H_j(\Gamma_N(m);\F)$ in the stated
degrees. Conjugating by permutation matrices gives the same surjectivity
on any specified $N-2$ coordinates, since each such matrix normalizes
$\Gamma_N(m)$ and therefore induces an automorphism of its homology.

Every elementary matrix $e_{ab}(f)$ centralizes the congruence subgroup
on the coordinates complementary to $a,b$. Its conjugation action
therefore fixes the image of a surjection onto $H_j(\Gamma_N(m);\F)$,
so that action is trivial. Elementary matrices generate $\SL_N(R)$
because $R$ is Euclidean, proving the lemma.
\end{proof}

We can now recover all the required congruence-group homology from its
coinvariants under $G_m$. The odd order of $G_m$ makes taking these
coinvariants exact, and the preceding lemma makes their action trivial.

\begin{proposition}\label{prop:quotient-homology}
For every $m\ge2$,
\[
 H_j(\Gamma_N(m);\F)=H_j(T_m;\F)=H^j(T_m;\F)=0
 \qquad(1\le j\le d-1).
\]
\end{proposition}

\begin{proof}
The finite group $G_m=U/\Gamma_N(m)$ has odd order: its quotient modulo
$t$ is the upper unitriangular group over $k$, and the kernel has a
filtration whose factors embed into additive matrix groups over $k$.
Let $P_*$ be a projective $\F[U]$-resolution of the trivial module $\F$.
Its restriction to $\Gamma_N(m)$ is still projective, because $\F[U]$
is free over $\F[\Gamma_N(m)]$. Hence
$(P_*)_{\Gamma_N(m)}$ computes $H_*(\Gamma_N(m);\F)$.
The quotient $G_m$ acts on this chain complex, and its action on homology
is the usual conjugation action. Averaging over the odd-order group
$G_m$ makes its coinvariant functor exact. Since
$P_U=(P_{\Gamma_N(m)})_{G_m}$, taking homology gives
\[
 H_j(U;\F)\cong H_j(\Gamma_N(m);\F)_{G_m}.
\]
By \cref{lem:trivial-action}, the right side equals
$H_j(\Gamma_N(m);\F)$ for $1\le j\le d-1$.
Both homology assertions now follow from
\cref{prop:homology-comparisons}, since $d-1<n-1$.
The cohomology assertion follows by duality over $\F$.
\end{proof}
